\documentclass[12pt, letterpaper]{report}% "letterpaper" matches US standard paper size.
\setcounter{tocdepth}{3}
% --- PACKAGES: Each one adds a capability to LaTeX ---
% Think of packages like Python imports or C #includes.


% ---------- Page Layout ----------
\usepackage[top=1in,bottom=1in,left=1in,right=1in]{geometry} % geometry lets you set margins. 1-inch all around is standard for reports.
% ---------- Encoding & Fonts ----------
\usepackage[utf8]{inputenc}% Allows special characters (é, ü, °) to be typed directly.
\usepackage[T1]{fontenc} % Ensures proper font encoding in the output PDF. Without this, copy-pasting from the PDF can produce garbled characters.
% ---------- Graphics & Figures ----------
\usepackage{graphicx}% THE core package for inserting images.
% Usage: \includegraphics[width=0.8\textwidth]{filename.png}
\graphicspath{{figures3/}}
\usepackage{float}% Gives you the [H] placement option for figures/tables. [H] means "put it HERE, not where LaTeX thinks is best." Very useful when you need a figure right after the text that references it.
\usepackage{subcaption}% Lets you put multiple images side-by-side in a single figure environment. Example: Figure 3(a) and Figure 3(b) next to each other.
% ---------- Tables ----------
\usepackage{booktabs} % Provides \toprule, \midrule, \bottomrule for professional-looking tables.
% Rule of thumb: NEVER use vertical lines in tables. Use booktabs rules instead.
\usepackage{tabularx}% Adds the tabularx environment, which lets columns auto-expand to fill width. Great for tables that need to span the full page width.
\usepackage{multirow}% Allows cells that span multiple rows in tables.
\usepackage{longtable}% For tables that are too long to fit on one page (like a Bill of Materials).
% ---------- Math ----------
\usepackage{amsmath}% The standard package for equations. Adds align, equation*, etc.
\usepackage{amssymb}% Extra math symbols (e.g., \therefore, \because, \pm).
\usepackage{siunitx}% Proper formatting for SI units.
\sisetup{per-mode=fraction}
\sisetup{exponent-product=\cdot}
\DeclareSIUnit{\fahrenheit}{\degree F}
% Usage: \SI{100}{\celsius}, \SI{10000}{psi}, \SI{10}{\centi\meter}
% This ensures consistent unit formatting throughout the document.
% ---------- Bibliography ----------
\usepackage[numbers, sort&compress]{natbib}
% IEEE-style numeric citations: [1], [2,3], [4-7]
% "sort&compress" turns [3,1,2] into [1-3]
% ---------- Cross-Referencing & Links ----------
\usepackage[
	colorlinks=true,
	linkcolor=blue,      % Internal links (sections, figures)
	citecolor=blue,      % Citation links
	urlcolor=blue         % URL links
]{hyperref}
% Makes all references, citations, and URLs clickable in the PDF.
% Set colorlinks=false if you want boxes instead of colored text.
% ---------- Captions ----------
\usepackage[
	font=small,           % Caption text slightly smaller than body
	labelfont=bf,         % "Figure 1:" is bold
	textfont=it            % The caption description is italic
]{caption}% Makes figure/table captions look more professional.
% ---------- Headers & Footers ----------
\usepackage{fancyhdr}% Customizable headers and footers.
\pagestyle{fancy}% Activate the fancy header/footer style.
\fancyhf{}% Clear all default header/footer content.
\fancyhead[L]{\small Team 16 --- Noble Aggie Labs}% Left header: team name.
\fancyhead[R]{\small EME 185}% Right header: course number.
\fancyfoot[C]{\thepage}% Center footer: page number.
\renewcommand{\headrulewidth}{0.4pt}% Thin line under the header. Set to 0pt to remove.
% ---------- Paragraph Formatting ----------
\usepackage{parskip}% Adds space between paragraphs instead of indenting the first line. This is more common in engineering reports than academic papers.
\setlength{\parindent}{0pt}% Explicitly remove paragraph indentation (parskip should handle this, but belt-and-suspenders doesn't hurt for readability).
% ---------- Line Spacing ----------
\usepackage{setspace}
% \onehalfspacing% 1.5 line spacing is easier to read and mark up than single spacing. Change to \doublespacing if required.
\onehalfspacing
%-------- For Box diagrams
\usepackage{tikz}
\usetikzlibrary{patterns}
\usetikzlibrary{patterns,intersections}
\usetikzlibrary{arrows.meta, positioning}
\usetikzlibrary{shapes.geometric}
\usetikzlibrary{arrows.meta, shapes.geometric, shapes.misc, positioning}
\usetikzlibrary{calc}
% ---------- Appendix Support ----------
\usepackage[title,titletoc,page]{appendix} % Adds "Appendices" header and includes appendices in the Table of Contents. [title, titletoc] ensures they are labeled A, B, C
%----------- Title/Chapter formatting -----------------
\usepackage{titlesec}
% Format \chapter to look simple (no giant text, no separate page)
\usepackage{pdfpages}
% Allows you to insert entire PDF Pages (e.g., for Gantt charts, concept screening matrices) with \includepdf.
\usepackage{geometry}
%allows you to set margins and page size. 
\usepackage{mdframed}     % warning/note boxes, important equations
	\mdfsetup{linewidth=1pt}
\usepackage{enumitem}     % list formatting
\usepackage{booktabs}     % tables
\usepackage{xcolor}       % color support
\usepackage{upgreek}


\definecolor{warnred}{RGB}{180,30,30}
\definecolor{warnbg}{RGB}{255,235,235}
\definecolor{notebg}{RGB}{235,245,255}
\definecolor{noteblue}{RGB}{30,80,160}
\definecolor{ucdgold}{RGB}{207, 184, 124}
\definecolor{ucdblue}{RGB}{0, 40, 85}
\definecolor{sandiablue}{RGB}{0,  174, 214}



\newmdenv[
  linecolor=warnred, backgroundcolor=warnbg,
  linewidth=1.5pt, innerleftmargin=10pt,
  innerrightmargin=10pt, innertopmargin=8pt,
  innerbottommargin=8pt, skipabove=10pt, skipbelow=6pt
]{warningbox}

\newmdenv[
  linecolor=noteblue, backgroundcolor=notebg,
  linewidth=1pt, innerleftmargin=10pt,
  innerrightmargin=10pt, innertopmargin=8pt,
  innerbottommargin=8pt, skipabove=10pt, skipbelow=6pt
]{notebox}


\titleformat{\chapter}
 {\LARGE\bfseries}  % format
 {}                % label
 {0em}                         % sep between label and title
 {}                            % before-code
\titlespacing*{\chapter}{0pt}{20pt}{10pt}  % left, before, after


\makeatletter
\g@addto@macro\appendix{%
 \titleformat{\chapter}[block]
	{\huge\bfseries\centering}
	{\appendixname\ \thechapter}
	{0em}
	{\vspace{0.5em}\\}%
}
\makeatother

% ---------- PDF Metadata ----------
\hypersetup{
    pdftitle={Final Design Report -- Passive Temperature Compensation Device},
    pdfauthor={Team 16: Noble Aggie Labs},
    pdfsubject={EME-185A/B Senior Design, UC Davis}
}
%-----------For landscape pages----------
\usepackage{pdflscape}
\usepackage{pgfplots}
	\pgfplotsset{compat=1.18}
	\usepgfplotslibrary{fillbetween}
	\pgfmathdeclarefunction{normalpdf}{3}{%
    \pgfmathparse{1/(#3*sqrt(2*pi))*exp(-(#1-#2)^2/(2*#3^2))}%
	}
\usepackage{multirow}
\usepackage{array}      % for >{\raggedright\arraybackslash}
\usepackage{makecell}   % for \makecell
\usepackage{calc}
\usepackage{multicol}
\usepackage{wrapfig}
\usepackage{nicefrac}
\usepackage[bottom]{footmisc}

% In your preamble
\newcounter{mycounter}           % creates a counter, initialized to 0
\newcounter{sub}[mycounter]      % resets 'sub' when 'mycounter' increments
% ============================================================================
% SECTION 0.5: CUSTOM COMMANDS (Modular Helpers)
% ============================================================================
% These are reusable shortcuts. Define them once, use them everywhere.
% If a name or term changes, you only edit it in ONE place.


% --- Project-specific shortcuts ---
\newcommand{\projecttitle}{Passive Temperature Compensation Device}
\newcommand{\teamname}{Noble Aggie Labs}
\newcommand{\teamnumber}{16}
\newcommand{\sponsor}{Dr.\ Bradley Bon}
\newcommand{\organization}{Sandia National Laboratories}
\newcommand{\coursename}{EME 185}
\newcommand{\instructor}{Professor Jonathon Schofield}
\newcommand{\university}{University of California, Davis}
\newcommand{\ta}{Teaching Assistant Xiangpu Wang}
\newcommand{\duedate}{5 June 2026}
% --- Convenience commands for common terms ---
% Using these ensures consistent terminology throughout the report.
\newcommand{\device}{PTCD}
\newcommand{\degC}{^{\circ}\text{C}}    % Degree Celsius symbol, properly formatted
\newcommand{\requiredRelationship}{D_{\text{effective}} = A_1 (T_{\text{ambient}} - T_{\text{reference}}) + A_0}
\newcommand{\transferFunction}{D_{\text{eff}}=\frac{dD_{\text{eff}}}{dX} \left[\frac{dL}{dT}\frac{dX}{dL}(T-T_{\text{ref}})+X\left(T_{\text{ref}} \right) \right]}

\newcommand{\Dcoldval}{}
\newcommand{\Dcold}{D_{\text{effective}}(\SI{-60}{\celsius}) = 383.5 \pm \SI{19.2}{\micro\meter}}
\newcommand{\Dhot}{D_{\text{effective}}(\SI{80}{\celsius}) = 170.2 \pm \SI{16.9}{\micro\meter}}
\newcommand{\Dref}{D_{\text{effective}}(\SI{25}{\celsius}) = 254.0 \pm \SI{12.7}{\micro\meter}}

\newcommand{\Thot}{\SI{80}{\degreeCelsius}} %Celcius
\newcommand{\Tcold}{\SI{-60}{\degreeCelsius}} %Celcius
\newcommand{\Tref}{\SI{25}{\degreeCelsius}} %Celcius
\newcommand{\specTempChange}{\SI{140}{\degreeCelsius}}

\newcommand{\specEAl}{\SI{68.95}{\giga \pascal}}
\newcommand{\specEALLV}{\SI{55}{\giga \pascal}}
\newcommand{\specEInv}{\SI{141}{\giga \pascal}}
\newcommand{\specEInc}{\SI{200}{\giga \pascal}}
\newcommand{\specActuatorCTE}{\SI{22.85}{\micro\meter\per\meter\per\celsius}} %Evaluated at 25C
\newcommand{\specMountCTE}{\SI{-30}{\micro\meter\per\meter\per\celsius}}
\newcommand{\specInvarCTE}{\SI{1.3}{\micro\meter\per\meter\per\celsius}}
\newcommand{\specInconelCTE}{\SI{13.2}{\micro\meter\per\meter\per\celsius}}
\newcommand{\specBellowsK}{\SI{75.3}{\kilo \newton \per \meter}} % 430 lb/in
\newcommand{\specBellowsHeight}{\SI{22.86}{\milli \meter}} % = Neck 1 length + Neck 2 length + Conv free diameter = 2*(0.07") + 0.76"
\newcommand{\specBellowsOD}{\SI{6.528}{\milli \meter}} % Conv OD 0.257"
\newcommand{\specBellowsID}{\SI{3.353}{\milli \meter}} % Conv ID 0.132"
\newcommand{\specBellowsSquirm}{\SI{20.68}{\mega \pascal}} % 3000 psi
\newcommand{\specBellowsBurst}{\SI{74.12}{\mega \pascal}} % 10750 psi

\newcommand{\Xhot}{\SI{480.696}{\micro \meter}}
\newcommand{\Xcold}{\SI{213.289}{\micro \meter}}
\newcommand{\specLengthChange}{\SI{267.407}{\micro \meter}}
\newcommand{\govErrorProp}{u_x = \sqrt{\frac{\pi(\sigma_x^2+\sigma_y^2)}{2}}}
\newcommand{\govActuator}{\Delta L = L_0 \Delta T \alpha}
\newcommand{\govValve}{D_{\text{eff}} = \sqrt{\frac{2}{\pi}}}
\newcommand{\specPressure}{\SI{68.95}{\mega \pascal}}%  68.94757 MPa
\newcommand{\AoneExact}{\SI{-1.524e-3}{\milli\meter\per\celsius}}
\newcommand{\Aone}{\AoneExact \pm 5\%}
\newcommand{\AzeroExact}{\SI{0.254}{\milli\meter}}
\newcommand{\Azero}{\AzeroExact \pm 5\%}
\newcommand{\imageHeight}{4in} % Default image height for figures. Adjust as needed.
\newcommand{\specCrossSectArea}{\SI{6.3e-5}{\meter \squared}}
\newcommand{\specOrigActuatorLength}{\SI{32.991}{\milli \meter}}
\newcommand{\specActDefTotal}{\SI{267}{\micro \meter}}
\newcommand{\specActDefHot}{\SI{162.4}{\micro \meter}}

% Transfer function specifications

\newcommand{\GA}{\SI{1.91}{\micro\meter\per\celsius}}
\newcommand{\GV}{\num{-0.798}}

\newcommand{\ReqA}{The device shall passively restrict gas flow according to the Required Relationship $D = A_1 \Delta T + A_0$}
\newcommand{\ReqB}{The device should remain within a \SI{10}{\centi\meter} $\times$ \SI{10}{\centi\meter} $\times$ \SI{10}{\centi\meter} mechanical envelope}
\newcommand{\ReqC}{The device shall restrict hydrogen at \specPressure\ and $T_{amb} + 20 \degC$ for 60 seconds in compliance with requirement R.1}
\newcommand{\RRfull}
{
	\begin{equation}
		\begin{aligned}
			D_{eff} &= \hat{G}_{V_0} \hat{G}_{A_0} \Delta T + D_{ref} \\
			&= \frac{dD}{dL} \frac{dL}{dT} (T_{amb} - T_{ref}) + D(T_{ref}) \\
			&= A_1 \Delta T + A_0
		\end{aligned}
	\end{equation}
}

% Valve specifications
\newcommand{\SpecVAA}{The valve subsystem shall passively control the effective orifice diameter as directly driven by the actuator at a conversion rate of $\frac{dD}{dL} = \GV$}
\newcommand{\SpecVAB}{The valve subsystem shall define the reference effective orifice diameter}
\newcommand{\SpecVABapp}{$D_{eff}(T_{ref}) = \AzeroExact$}
\newcommand{\SpecVB}{The valve subsystem shall not exceed a \SI{10}{\centi\meter} $\times$ \SI{5}{\centi\meter} $\times$ \SI{10}{\centi\meter} mechanical envelope}
\newcommand{\SpecVCA}{The housing shall withstand an internal pressure of \specPressure}
\newcommand{\SpecVCB}{Wetted components shall maintain a hermetic seal at this pressure}
\newcommand{\SpecVCC}{Wetted components shall not permit hydrogen diffusion or embrittlement}

% Actuator specifications
\newcommand{\SpecAAA}{The actuator subsystem shall passively drive the valve subsystem in response to ambient temperature at a rate of $\frac{dL}{dT} = \GA$}
\newcommand{\SpecAAB}{The actuator subsystem shall facilitate user calibration of the device}
\newcommand{\SpecAB}{The actuator subsystem shall not exceed a \SI{10}{\centi\meter} $\times$ \SI{5}{\centi\meter} $\times$ \SI{10}{\centi\meter} mechanical envelope}
\newcommand{\SpecACA}{The actuator shall allow no more than \textbf{\textcolor{red}{[TODO: Spec Temp]}} deviation in its temperature from $T_{amb}$\ within the test window}

% --- Figure insertion helper ---
% This command simplifies inserting a figure with a caption and label.
% Usage: \insertfigure{width}{filename}{caption text}{fig:label}
%
% Parameters:
%   #1 = width relative to text width (e.g., 0.8 for 80%)
%   #2 = filename (without path if images are in the same directory)
%   #3 = caption text
%   #4 = label for cross-referencing (use fig:prefix by convention)
\newcommand{\insertfigure}[4]{%
	\begin{figure}[H]
		\centering
		\includegraphics[width=#1\textwidth]{#2}
		\caption{#3}
		\label{#4}
	\end{figure}
}

% --- TODO marker for draft mode ---
% Prints a bold red "TODO" so you can easily spot incomplete sections.
\usepackage{xcolor}
\usepackage{tocloft} % For listoftodos
% --- TODO marker ---
\newlistof{todos}{tod}{List of TODO's}
\newcommand{\todo}[1]{%
    \textbf{\textcolor{red}{[TODO: #1]}}%
    \addcontentsline{tod}{todos}{\protect\numberline{\thepage}#1}%
}

% For links to spec table
\newcommand{\specref}[1]{\hyperref[tab:specs]{#1}}

\usepackage{lmodern} % DELETE BEFORE SUBMISSION; this is to help the pdf render more easily.

\begin{document}
\chapter{Testing}
\label{chap:test}

\todo{Outline for Testing chapter. Combine isolated Heat Transfer section with Miguel's. Obejective = SAME, Test Design = SAME, Boundary Conditions = SAME, Defomation and Dimensional Tolerance Flow Chart = SAME, 
Mesh Decisions = DIFFERENT, DATA COLLECTION = DIFFERENT. The "DIFFERENT" sections can still be married and ideas synthesized. We need to include limitations and a disucssion for results. Can reference appendix if needed for space. 
We need to split limitations in subsections, for each test CFD, MBD, and HT. For heat transfer limitation, contact resistance was omitted in these tests. Limitation in MBD where geometry was not decoupled at the orifice, assumes parts are welded 
together, MBD different in that inlet and outlet were omitted, but pressure left in from CFD. Discuss but do not synthesize for all limitations!}
\section{Objective}
In the absence of a physical manufactured model, verification of the DFMEA is done through two methods: (1)analytical methods, (2) Finite element simulations. This section discusses the testing process and strategy of using finite element simulations to validate the analytical results and explore failure modes determined by DFMEA.  The testing plan validated and expanded/iterated on analytical methods by simulating the entire PTCD. Testing includes conducting Heat Transfer, Computational Fluid Dynamics (CFD), and Multibody Dynamics (MBD) in isolation, as well as combined.

\section{Test Design}

\todo{Add corresponding UI screenshots of COMSOL that reflect the boundary condition decisions. Include the COMSOL tree and the actual corresponding part of the model.}

\subsection{Set-up: General Strategy}

To make the simulations computationally less expensive, the following strategies have been employed
\begin{enumerate}
	\item The PTCD is cut in half about its vertical symmetry plane, creating new surfaces which will be referred to as "symmetry surfaces". By setting the exposed symmetry surfaces to its appropriate "symmetrical nodes" in the simulation software, the amount of meshes can be cut in half, reducing the computation time. \todo{INSERT SYMMETRY FACES PHOTOS HERE}
	\item The physics is "decoupled". The CFD is ran on a rigid negative domain of the channel with a turbulent simulation where a "pressure map" of the wetted surface is recovered. This same pressure map is used to apply a load on the wetted surfaces of the PTCD assembly geometry to apply a load that is representative of the fluid flow during operation. The heat transfer simulation is then executed on the loaded assembly by setting all exposed surfaces to the ambient temperature, and all the wetted surfaces to the fluid temperature.
\end{enumerate}

\subsection{Set-up: Boundary Condition Decisions}
\begin{itemize}
    \item Surface pairs and contact constraints.
	\\
	In multi-body simulation, the consideration for surface pairs where the multiple bodies contact is pivotal for accuracy of results. For all simulations, the welded surfaces are set as \textbf{identity pair} while surfaces that are unwelded and expected to collide are set as \textbf{contact pairs}.
	Failure to set up these surface constraints will result in bodies phasing through each other, welded surfaces to decouple, and so on.
	\\
	Verification for the correctness of contact constraints are made by setting a rigid translation of the orifice in space and observing how the meshes behave. This is shown on the \todo{insert figure reference}figure below

    \item Fluid-solid boundaries (incorporating union nodes where a single continuous boundary is required for FSI).
    \item Accurate boundary representations for complex orifice shapes, ensuring internal flow geometries are not oversimplified as standard cylinders.
\end{itemize}

\subsection{Deformation and Dimensional Tolerance Flow Chart}
Mapping the transfer function dependencies: CFD $\rightarrow$ MBD $\rightarrow$ Deformation $\leftarrow$ Heat Transfer.

\subsection{Mesh Decisions}
Strategies and refinements to improve computational accuracy and validity across domains.

\section{Data Collection}

\subsection{Design Space Matrix}
\begin{itemize}
    \item Turbulent flow $\Delta P$ range.
    \item Heat transfer: $T_{amb}$ range and $T_{gas}$ range.
\end{itemize}

\section{Results}

\section{Limitations}

\chapter{Testing}
\label{chap:test}

\todo{Outline for Testing chapter}
\section{Objective}
Validate and expand/iterate on analytical methods by simulating the entire PTCD. Testing is purely simulation-based in COMSOL Multiphysics and includes conducting Heat Transfer (thermal-stress) and transient thermal studies, both for the actuator in isolation and for the full core assembly.

\section{Test Design}

\subsection{Set-up: Boundary Condition Decisions}
\begin{itemize}
    \item Surface pairs and contact constraints (Form Assembly with Identity Pairs for the multi-body assembly; Rigid Connector face bonding for the simplified actuator–mount model).
    \item Bellows represented as an equivalent cylinder with softened effective modulus to match axial spring stiffness, avoiding intractable bellows-geometry meshing.
    \item Thermal environment applied through the $T_{env}$ parameter, with a fixed constraint at the mount base providing a stable reference for actuator displacement.
\end{itemize}

\subsection{Stroke and Dimensional Tolerance Flow Chart}
Mapping the transfer function dependencies: $T_{env} \rightarrow$ Heat Transfer $\rightarrow$ Thermal Expansion $\rightarrow$ Actuator Stroke $\rightarrow$ Dimensional Tolerance, with elastic deformation of the mount and actuator as a parallel contribution.

\subsection{Mesh Decisions}
Strategies and refinements to improve computational accuracy and validity (physics-controlled meshing, refinement at flexures and junctions, isolation of sharp-corner singularities from bulk results).

\section{Data Collection}

\subsection{Design Space Matrix}
\begin{itemize}
    \item Heat transfer: $T_{amb}$ range ($-60\,^{\circ}\mathrm{C}$ to $+80\,^{\circ}\mathrm{C}$) swept about the $25\,^{\circ}\mathrm{C}$ reference.
    \item Mount material candidates: Invar 36 and Allvar Alloy 30 (differential-CTE comparison).
    \item Transient response: settling time to steady state for a representative $\Delta T$.
\end{itemize}

\subsection{Validation Against Analytical Model}
\begin{itemize}
    \item Simulated actuator displacement (point evaluation on the actuator bottom face) overlaid against the analytical linear model.
    \item Percent error between FEA and analytical stroke across the operating range.
\end{itemize}

\section{Data Analysis}

\subsection{Dimensional Failure}
Focusing on heat transfer impacts and thermal expansion, including verification of the linear stroke relationship and the behavior of wetted surfaces and core-assembly components across the operating range.

\subsection{Stress and Material Failure}
\todo{Show images of the failure modes reflected on the model. An optional suggestion to have some sort of validation reference via a figure from a study, etc.}

Evaluating structural integrity, yielding, and mechanical limits of the actuator, mount, valve, and bellows under thermal load.

\end{document}